\documentclass{article}
\usepackage[utf8]{inputenc}
\usepackage{geometry}
\usepackage{verbatim} % Needed for the code block
\geometry{a4paper, margin=1in}

\title{Análisis del Proyecto de Seguridad Informática}
\author{José}
\date{\today}

\begin{document}

\maketitle

\section{Análisis del Trabajo Asignado}
La tarea asignada a José consistía en implementar la función \texttt{scan\_ports(target, port\_spec)} dentro del archivo \texttt{modulos.py} para añadir la funcionalidad de escaneo de puertos a la herramienta.

\subsection{Resultado de la Tarea}
La tarea se encuentra \textbf{incompleta}. La función \texttt{scan\_ports} fue definida en el archivo, pero no contiene ninguna lógica para realizar el escaneo. En su lugar, simplemente imprime un mensaje indicando que el módulo no está implementado.

El código actual de la función es el siguiente:
\begin{verbatim}
def scan_ports ( target , port_spec ) :
    """
    Escanea los puertos especificados en el host objetivo .
    ( Implementacion pendiente para futuras tareas )
    """
    print ( f"[!] El modulo de escaneo de puertos para '{ port_spec }' aunesta implementado .")
    return []
\end{verbatim}

\section{Estructura del Proyecto}
Se observó que José no trabajó sobre la base de código común del proyecto. En su lugar, creó una estructura de proyecto paralela que incluye su propio script principal (\texttt{Auditoria.py}) y un archivo de módulos (\texttt{modulos.py}) que duplica código de otros miembros del equipo (como las funciones de consulta DNS).

Esta aproximación sugiere una falta de sincronización y colaboración con el resto del equipo, lo que dificulta la integración del trabajo en un único proyecto funcional.

\section{Conclusión}
El trabajo asignado a José \textbf{no fue completado} y la funcionalidad principal que debía desarrollar está ausente. Además, la duplicación de la estructura del proyecto indica un problema de coordinación que debería ser atendido.

\end{document}
